\section*{Capítulo \ref{ch:b3:spin-two-level} --- Spin e sistemas de dois níveis}

\begin{solution}{exo:b3:spin-two-level:1}
(a) Multiplicação direta. (b) $[\hat S_x, \hat S_y] =
(\hbar/2)^2\,2\iu\sigma_z = \iu\hbar\hat S_z$. (c) $\hat S^2 =
(\hbar/2)^2 \times 3\,\mathbb 1 = \tfrac34\hbar^2$: de fato,
$\tfrac12\cdot\tfrac32\hbar^2$. (d) $\{\mathbb 1, \sigma_x,
\sigma_y, \sigma_z\}$ geram as quatro dimensões reais das matrizes
hermitianas $2\times2$ ($a + \vect b\cdot\vect\sigma$ com $a,
\vect b$ reais).
\end{solution}

\begin{solution}{exo:b3:spin-two-level:2}
(a) $\ket{\pm x} = (\ket\uparrow \pm \ket\downarrow)/\sqrt2$;
$\ket{\pm y} = (\ket\uparrow \pm \iu\ket\downarrow)/\sqrt2$. (b)
Metade e metade. (c) Metade e metade de novo: $|\braket{\pm y}{+x}|^2 =
|1 \mp \iu|^2/4 = \tfrac12$. (d) $\ket{+x}$: equador em $\varphi =
0$; medir ao longo de $z$ ou de $y$ projeta sobre polos a $90^\circ$
--- sempre $\cos^2(45^\circ) = \tfrac12$.
\end{solution}

\begin{solution}{exo:b3:spin-two-level:3}
(a) $\cos^2(\theta/2)$. (b) $\cos^2(\theta/2) \times
\cos^2(\theta/2) = \cos^4(\theta/2)$. (c) $\theta = 90^\circ$:
$(\tfrac12)^2 = \tfrac14$, o número da cadeia de polarizadores. (d)
Trivialmente $\theta = 0$; a lição de verdade é que as medidas
intermediárias \emph{suaves} custam menos --- o limite de passos
pequenos do \cref{exo:b3:quantum-formalism:5}, em que muitas pequenas
rotações passam o estado adiante com perda evanescente.
\end{solution}

\begin{solution}{exo:b3:spin-two-level:4}
(a) $2\mu_{\text{B}}B = \qty{116}{\micro eV}$: $\nu =
\qty{28}{GHz}$. (b) $2 \times 2.79\,\mu_{\text{N}}B =
\qty{0.18}{\micro eV}$: $\nu = \qty{42.6}{MHz}$. (c) Contra
$k_{\text{B}}T = \qty{25.9}{meV}$: polarização eletrônica de
$\sim\num{2e-3}$ e protônica de $\sim\num{3e-6}$ --- os conjuntos
térmicos de \omterm{def:b3:spin-two-level:spin}{spins} estão quase perfeitamente embaralhados. (d) A RPE:
micro-ondas; a RMN: rádio --- três ordens de grandeza de distância,
uma só física.
\end{solution}

\begin{solution}{exo:b3:spin-two-level:5}
(a) $\dd\langle\hat S_x\rangle/\dd t = -\gamma B_0\langle\hat
S_y\rangle$, $\dd\langle\hat S_y\rangle/\dd t = \gamma
B_0\langle\hat S_x\rangle$, $\dd\langle\hat S_z\rangle/\dd t = 0$.
(b) $\langle\hat S_x\rangle = \tfrac\hbar2\cos\omega_0t$,
$\langle\hat S_y\rangle = \tfrac\hbar2\sin\omega_0t$ (com o sinal
conforme o de $\gamma$). (c) As duas coisas são evidentes em (b). (d) O
estado é uma \emph{superposição} dos dois autoestados de energia; sua
fase relativa avança a $\omega_0$, e essa fase que gira \emph{é} a
precessão --- níveis estacionários e vetor que gira são as duas
leituras de uma só evolução de dois níveis.
\end{solution}

\begin{solution}{exo:b3:spin-two-level:6}
(a) Substitua e use as identidades de meio ângulo. (b)
$\cos\tfrac\theta2\cos\tfrac{\pi-\theta}2 +
\eu^{\iu\pi}\sin\tfrac\theta2\sin\tfrac{\pi-\theta}2 =
\cos\tfrac\theta2\sin\tfrac\theta2 -
\sin\tfrac\theta2\cos\tfrac\theta2 = 0$. (c) No equador, em
$\varphi = 0, \pi$ ($\sigma_x$) e $\varphi = \pm\pi/2$
($\sigma_y$). (d) Uma rotação de $180^\circ$ em torno de qualquer eixo
equatorial --- o pulso $\pi$ da \omterm{prop:b3:spin-two-level:rabi}{ressonância magnética} é a porta NÃO do
qubit.
\end{solution}

\begin{solution}{exo:b3:spin-two-level:7}
(a) No referencial girante a excitação fica parada, ao passo que a
rotação do referencial subtrai $\omega/\gamma$ do campo estático
efetivo (a mesma contabilidade de uma força de inércia em referencial
girante). (b) Só $B_1$ sobrevive: o vetor de Bloch precessa em torno do
eixo transversal (fixo) de $B_1$ a $\Omega_1 = |\gamma|B_1$ --- viradas
regulares. (c) $t_\pi = \pi/\gamma B_1 = 1/(2 \times 42.58\,\text{
MHz/T} \times 10\,\unit{\micro T}) = \qty{1.2}{ms}$. (d) Uma espiral
rápida que fecha o cone: precessão a \qty{128}{MHz} em torno de $z$,
com nutação lenta descendente a \qty{426}{Hz} --- um saca-rolhas de mil
voltas, de polo a polo.
\end{solution}

\begin{solution}{exo:b3:spin-two-level:8}
(a) $\nu = \Delta E/h = \num{5.9e-6}/\num{4.14e-15} =
\qty{1.42e9}{Hz}$; $\lambda = c/\nu = \qty{21.1}{cm}$. (b) $\Delta
\nu \sim 1/2\pi\tau \sim \qty{e-15}{Hz}$: completamente desprezível ---
toda largura observada é movimento. (c) $\Delta\nu =
\nu\,(2v/c) = 1420\,\text{MHz} \times \num{1.5e-3} \approx
\qty{2}{MHz}$: o perfil de dois chifres de um disco em rotação. (d)
A de 21 cm fala onde o hidrogênio é atômico (morno, difuso); o CO fala
onde o hidrogênio se emparelhou no invisível H$_2$ (frio, denso):
juntos, eles inventariam todo o meio interestelar de uma galáxia.
\end{solution}

\begin{solution}{exo:b3:spin-two-level:9}
(a) Duas configurações acopladas por tunelamento: em baixa energia o
espaço é bidimensional --- a amônia \emph{é} um \omterm{def:b3:spin-two-level:spin}{spin} $\tfrac12$
disfarçado. (b) $\nu = \Delta E/h = \qty{24}{GHz}$, $\lambda =
\qty{1.25}{cm}$. (c) O feixe precisa chegar com o estado
\emph{superior} superpovoado (selecionado por deflexão eletrostática)
--- inversão de população, inatingível termicamente, pois Boltzmann
sempre favorece o nível de baixo. (d) A amplificação estimulada por
inversão de população --- o princípio de funcionamento do laser,
demonstrado primeiro em frequências de micro-ondas.
\end{solution}

\begin{solution}{exo:b3:spin-two-level:10}
(a) $\mu_{\text{B}} \times \qty{10}{T} = \qty{0.58}{meV}$ --- a escala
$\alpha^2E_{\text{I}} \approx \qty{0.7}{meV}$ do
\cref{exo:b3:hydrogen-atom:11}: a estrutura fina \emph{é} o \omterm{def:b3:spin-two-level:spin}{spin}
encontrando o campo do movimento. (b) $\Delta E = hc\,\Delta\lambda/
\lambda^2 = \qty{2.1}{meV}$: um campo interno $\Delta E/2\mu_{
\text{B}} \approx \qty{18}{T}$. (c) O campo interno cresce como
$Z^4$ (a carga nuclear ao cubo no campo, e mais uma vez no raio da
órbita): os átomos pesados têm uma estrutura fina que se vê com um
espectroscópio de bolso. (d) $j = \tfrac32$: quatro estados;
$j = \tfrac12$: dois --- seis ao todo, exatamente $2 \times (2l + 1)$
para $l = 1$.
\end{solution}

\begin{solution}{exo:b3:spin-two-level:11}
(a) Contagens de $S_z$: $+\hbar, 0, 0, -\hbar$; agir com $\hat S^2$
(pela identidade da escada) sobre a combinação simétrica dá
$2\hbar^2$ ($S = 1$), e sobre a antissimétrica dá $0$ ($S = 0$). (b)
Só o singleto (e o estado do meio do tripleto) não pode ser escrito
como produto --- o singleto é \emph{o} par maximamente emaranhado. (c)
Emissão significa que o tripleto está acima: a configuração de \omterm{def:b3:spin-two-level:spin}{spins}
paralelos é a mais energética, por \qty{5.9}{\micro eV}. (d)
Cada observador, sozinho, vê resultados perfeitamente aleatórios; a
correlação só aparece quando as duas listas são \emph{reunidas} ---
nenhuma estatística local muda, e por isso nada se propaga
(o argumento da \cref{rem:b3:quantum-formalism:nosignal}).
\end{solution}

\begin{solution}{exo:b3:spin-two-level:12}
(a) O primeiro $\pi/2$ deita o vetor de Bloch no equador; durante
$T$ ele precessa no $\nu_0$ próprio do átomo, enquanto o oscilador
local gira a $\nu$: a diferença de ângulo acumulada vale
$2\pi\delta T$; o segundo $\pi/2$ converte essa fase numa população
--- um interferômetro no tempo. (b) $\Delta\delta \sim
1/2T = \qty{1}{Hz}$. (c) $\qty{e-4}{Hz}/\qty{9.19e9}{Hz} \approx
\num{e-14}$ --- um segundo em três milhões de anos. (d) A mesma divisão
absoluta de franja sobre uma portadora $10^5$ vezes mais alta ganha
$10^5$ em exatidão relativa: daí os relógios ópticos de estrôncio e de
itérbio em $10^{-18}$, e uma redefinição pendente do segundo.
\end{solution}

\begin{solution}{pb:b3:spin-two-level:1}
\textbf{1.} $\nu = 42.58 \times 3.0 = \qty{127.7}{MHz}$; $h\nu =
\qty{5.3e-7}{eV}$.
\textbf{2.} $\Delta n/n = h\nu/2k_{\text{B}}T = \num{5.3e-7}/(2
\times 0.0267) \approx \num{e-5}$: dez partes por milhão.
\textbf{3.} $\num{6.6e22}\,\unit{cm^{-3}} \times \num{e-5} \approx
\num{7e17}$ prótons portadores de sinal por centímetro cúbico ---
débeis por \omterm{def:b3:spin-two-level:spin}{spin}, poderosos em número.
\textbf{4.} Stern--Gerlach mediu cada átomo e o separou por \omterm{def:b3:quantum-formalism:observable}{autovalor};
a \omterm{prop:b3:spin-two-level:rabi}{ressonância magnética} escuta uma \emph{média térmica} de $10^{23}$
\omterm{def:b3:spin-two-level:spin}{spins}, e Boltzmann mantém essa média a microvolts do zero.
\textbf{5.} A polarização $\propto B_0$ e a fem induzida
$\propto \omega \propto B_0$: sinal aproximadamente $\propto B_0^2$ ---
o argumento a favor de \qty{3}{T} contra \qty{1.5}{T}, e das máquinas
de pesquisa de \qty{7}{T}.
\textbf{6.} Um ferromagneto no gradiente de fuga sente
$F = \mu\,\partial B/\partial z$ escalado a quilogramas de ferro:
centenas de newtons surgindo numa porta --- a razão da blindagem e das
listas de verificação.
\textbf{7.} $\nu_1 = \gamma B_1/2\pi = \qty{426}{Hz}$: pulso $\pi/2$ de
\qty{0.59}{ms}, pulso $\pi$ de \qty{1.2}{ms}.
\textbf{8.} Ela precessa no plano transversal a \qty{127.7}{MHz}; o
fluxo da magnetização que gira, através da bobina, induz (Faraday) uma
fem de rádio exatamente nessa frequência --- o sinal bruto da
ressonância.
\textbf{9.} Toda virada que troca energia também torna aleatória a
fase do \omterm{def:b3:spin-two-level:spin}{spin} virado: o que quer que cause $T_1$ contribui para $T_2$, e
o desfasamento ainda tem canais próprios --- a coerência morre
primeiro.
\textbf{10.} Em $\tau$ todo corredor dá meia-volta: os mais rápidos,
que estão mais à frente, têm agora mais caminho de volta; em $2\tau$
todos cruzam a linha de partida lado a lado --- o desfasamento se
rebobina e o eco soa.
\textbf{11.} Só a parte \emph{estática}: um \omterm{def:b3:spin-two-level:spin}{spin} que manteve uma
frequência constante (ainda que errada) refaz sua fase exatamente; os
campos moleculares flutuantes mudam entre as duas metades e não se
rebobinam --- o decaimento deles é o verdadeiro $T_2$, específico do
tecido.
\textbf{12.} Amostre cedo e com frequência (TR curto, TE curto) e o
$T_1$ curto da gordura brilha; espere muito e ecoe tarde e os líquidos
de $T_2$ longo é que reluzem: os ajustes de relógio da sequência
escolhem que constante de relaxação pinta a imagem.
\textbf{13.} $\dd\nu/\dd z = 42.58\,\unit{MHz/T} \times
0.04\,\unit{T/m} = \qty{1.7}{kHz/mm}$.
\textbf{14.} A fatia em que a frequência de Larmor local cai na banda:
espessura $\qty{2}{kHz}/\qty{1.7}{kHz/mm} \approx
\qty{1.2}{mm}$ --- uma fatia selecionada.
\textbf{15.} Uma transformada de Fourier: a frequência rotula a
posição, de modo que o espectro do eco \emph{é} a projeção da fatia.
\textbf{16.} $256^2 \times 12 \approx \qty{100}{kB}$; a uma linha
codificada por tempo de repetição da ordem de $T_1$, $256$ linhas
custam minutos --- por isso se pede aos pacientes que fiquem parados.
\textbf{17.} As bobinas de gradiente conduzem correntes chaveadas de
escala quiloampère dentro de \qty{3}{T}: a força de Laplace as martela
contra os suportes a cada chaveamento --- a trilha sonora de
metralhadora de todo exame.
\textbf{18.} Os raios X sombreiam a densidade eletrônica --- osso
contra ar --- e depositam dose ionizante; a ressonância lê a densidade
de prótons e dois relógios de relaxação, que diferem ricamente entre os
tecidos moles: cérebro, cartilagem e tumores, todos iguais para os
raios X, são distintos para os \omterm{def:b3:spin-two-level:spin}{spins}, com dose ionizante nula.
\textbf{19.} Cada ambiente químico blinda seu próton em algumas partes
por milhão: os prótons de uma molécula se apresentam como um pente
resolvido de linhas deslocadas --- determinação de estrutura num tubo.
\textbf{20.} A atividade aumenta o fluxo sanguíneo e a oxigenação
locais; a desoxi-hemoglobina paramagnética se dilui; o $T_2^*$ local se
alonga; o voxel clareia segundos depois do pensamento.
\textbf{21.} A \qty{84}{GHz} o tecido é opaco (as micro-ondas
milimétricas mal penetram a pele) e a relaxação eletrônica é de
microssegundos --- inútil in vivo, mas perfeita para contar radicais e
defeitos em materiais e em dosimetria.
\textbf{22.} $\num{5e-7}\,\unit{eV}$ contra ligações de eV: dez milhões
de vezes fraca demais para quebrar o que quer que seja --- não
ionizante; um único fóton de raios X de \qty{60}{keV} carrega
$10^{11}$ vezes mais.
\textbf{23.} De $10^{-5}$ à ordem de um: um ganho de $\sim10^5$ por
núcleo --- necessário porque o gás inalado é milhares de vezes mais
diluído que os prótons da água.
\textbf{24.} Um \omterm{def:b3:spin-two-level:spin}{spin} $\tfrac12$ é um sistema de dois níveis
\emph{perfeito} (sem níveis de fuga), e se acopla ao ruído de carga
apenas fracamente, por momentos magnéticos: coerência longa num
material industrial.
\textbf{25.} Uma precessão de \qty{128}{MHz}, uma polarização de
$10^{-5}$, dois relógios de relaxação e três gradientes: a física do
\omterm{def:b3:spin-two-level:spin}{spin} fotografando um cérebro vivo, fatia a fatia, com dose ionizante
nula --- as duas manchas de prata de 1922, crescidas até virar um
departamento de hospital.
\end{solution}
